\honest{Go Forth and Create the Art!}{Eliezer Yudkowsky}{2009}

I have said a thing or two about rationality: how to untangle confused questions, how to tell real reasoning from fake, how to try before you flee, how to do the impossible. I developed these techniques in my own projects, and your mileage may vary. Still, knowing the conjunction fallacy hardly seems esoteric. Understanding why motivated skepticism is bad for you may be the whole difference between a smart person who stays smart and one who ends up stupid, and affective death spirals consume many of the unwary.

There is more absent than present in this art. I have not solved akrasia or how to coordinate groups, I have worked more on belief than on action, and there is still no training, teaching or verification. Building the art could also mean better introductory writing, better slogans, common cause with other Enlightenment projects, and addressing the gender imbalance. \nb{``A Sense That More Is Possible'' and ``3 Levels of Rationality Verification,'' earlier in this sequence, made these points at length.}

I suspect, ``you could even call it a guess,'' that there is a barrier to getting started. Anyone who begins to sense that more is possible may go wrong at once. As David Stove observes, most great philosophers, Hegel for one, ``are properly objects of pity.'' Developing the art of thinking is not unlike what I was doing in AI: you are tempted by fake explanations of the mind, fake accounts of causality, mysterious holy words and the amazing idea that solves everything. People who start thinking about thinking go off and invent Freudian psychoanalysis, or a new religion. \nb{Stove's essay, ``What Is Wrong with Our Thoughts?'', says this: the ``great thinkers'' are ``proper objects, indeed, of pity, but even more, of horror,'' and people who attempt ``any depth or generality of thought'' go mad ``almost infallibly.'' Yudkowsky gave the essay its own post a few weeks later.}

My fake-detecting methods are not best for every problem, but they may help to tell good systems of thinking from bad. I hope that what I have set down gets you past that barrier: that you will not run away or make things up, that you may consult experimental psychology, will not spiral around your Brilliant Idea, will know a fake explanation from a real one, and that you ``will get a saving throw.'' That is my dream: that this specialized art of answering confused questions may be some of what is needed to complete the rest.

The rest is ``a task which I am leaving to you. Probably, anyway.'' I have other things to do. \nb{For information: the Center for Applied Rationality was founded in 2012 to teach rationality in workshops. Its own evaluation is reported in the notes on ``3 Levels of Rationality Verification.''}

Two pieces of advice. First, draw on many sources, not on one author, though perhaps on a central website to post the links and papers that matter. ``To the best of my knowledge'' only cults draw their strength from one person; a true science may have lonely defiant heroes, but more than one. Second, develop the art while trying to do something that matters, perhaps a task hard enough to break your old understanding, not by sitting and asking how to fight akrasia. Perhaps the next work will instead come from studying rationality directly. I hedge: ``Maybe I'm being too idealistic,'' that I may be ``overgeneralizing from my own history,'' and ``maybe I'm wrong.'' \nb{In ``Something to Protect'' (2008) Yudkowsky had put this without a hedge: ``No one masters the Way until more than their life is at stake.''}

My post on secret identities was not well received; ``leave the house'' may be better. I am not sure I ever put the central rhythm of the art into words; I can say the words but not the rule that generates them; perhaps using the ideas will grow similar machinery in you. All human efforts to learn arcana slide by default into passwords, hymns and floating assertions. My earlier efforts to teach it were mostly unsuccessful; this time I threw a larger rock, and time will tell. I hope the karma system will stop the decline that has ended every community I tried to build.

I want people to go forth, but also to return, or even to go and stay at once, since seeing that others are trying helps motivation over years. So go forth, confront challenges, create new art beyond your armchair, and ``radio back to tell others what you learned.''
